% ============================================================
% Comparison: WAAM-MultiLayer-Simulation (original) vs WAAM-Framework
% Compile: pdflatex ORIGINAL_vs_FRAMEWORK_COMPARISON.tex
% ============================================================
\documentclass[10pt,a4paper]{article}

\usepackage[top=2.0cm, bottom=2.0cm, left=1.8cm, right=1.8cm]{geometry}
\usepackage[utf8]{inputenc}
\usepackage{amsmath,amssymb,amsfonts}
\usepackage{mathptmx}
\usepackage{microtype}
\usepackage{booktabs}
\usepackage{tabularx}
\usepackage{array}
\usepackage{longtable}
\usepackage{multirow}
\usepackage{xcolor}
\usepackage{titlesec}
\usepackage{fancyhdr}
\usepackage{setspace}
\usepackage{enumitem}
\usepackage{ragged2e}
\usepackage[colorlinks=true,linkcolor=black,citecolor=black,urlcolor=black]{hyperref}

\definecolor{rowalt}{HTML}{F2F2F2}
\setlength{\extrarowheight}{1.5pt}
\renewcommand{\arraystretch}{1.15}

\setstretch{1.08}
\setlength{\parindent}{0pt}
\setlength{\parskip}{0.45em}

\pagestyle{fancy}
\fancyhf{}
\fancyhead[L]{\small\textit{Original vs.\ Framework --- WAAM Simulation Code Comparison}}
\fancyhead[R]{\small\textit{P. K. Pandey (SC25M147)}}
\fancyfoot[C]{\thepage}
\renewcommand{\headrulewidth}{0.4pt}
\renewcommand{\footrulewidth}{0pt}

\titleformat{\section}{\large\bfseries}{\thesection.}{0.7em}{}[\titlerule]
\titleformat{\subsection}{\normalsize\bfseries}{\thesubsection}{0.5em}{}
\titleformat{\subsubsection}{\small\bfseries\itshape}{\thesubsubsection}{0.4em}{}
\titlespacing*{\section}{0pt}{1.4em}{0.6em}
\titlespacing*{\subsection}{0pt}{1.1em}{0.4em}
\titlespacing*{\subsubsection}{0pt}{0.9em}{0.3em}

\newcommand{\yes}{\textcolor{blue!55!black}{\textbf{Yes}}}
\newcommand{\no}{\textcolor{red!60!black}{\textbf{No}}}
\newcommand{\na}{---}
\newcommand{\code}[1]{\texttt{#1}}
\newcommand{\file}[1]{\texttt{\small #1}}

\begin{document}

% ------------------------------------------------------------
\begin{center}
    {\LARGE\bfseries A Detailed Comparison of Two WAAM Simulation Codes\par}
    \vspace{0.5em}
    {\large\textbf{Original \textit{WAAM-MultiLayer-Simulation} (FreeFEM++ single
    script) versus the derived \textit{WAAM-Framework} (multiphysics FEM
    platform)}\par}
    \vspace{0.7em}
    {\large P.~K.~Pandey (SC25M147) \quad$\cdot$\quad IIST \quad$\cdot$\quad
    13 October 2026\par}
\end{center}
\vspace{0.3em}

\noindent\fbox{\parbox{\dimexpr\linewidth-2\fboxsep-2\fboxrule}{%
\small\textbf{Scope of this document.}
Two code bases held under \code{Desktop/WAAM/} were compared line by line:
(i) the \emph{original} upstream repository
\code{WAAM-Framework/original\_git\_file/} containing \file{README.md}
(434~lines) and \file{downloaded.edp} (534~lines) --- a single hand-written
FreeFEM++ script; and (ii) the \emph{framework} \code{WAAM-Framework/},
containing 28{,}595~lines of Python across 40~source files, 3{,}571~lines of
reference \code{.edp} modules, and a code generator that emits scripts of up to
$\sim$3{,}150~lines. All line counts, file names and line references below were
measured from the working tree on 13 October 2026.}}

\vspace{1.0em}

% ============================================================
\section{Executive Summary}
% ============================================================

\begin{table}[h!]
\centering
\caption{Head-to-head summary of the two code bases.}
\label{tab:summary}
\small
\begin{tabularx}{\linewidth}{@{}l X X@{}}
\toprule
\textbf{Dimension} & \textbf{Original} & \textbf{Framework} \\
\midrule
Nature & One hand-written FreeFEM++ script plus a README &
Layered Python application that \emph{generates} FreeFEM++ scripts \\
\addlinespace
Scale & 968 lines, 2 files & 28{,}595 lines Python + 3{,}571 lines
\code{.edp} + generator ($\sim$40$\times$) \\
\addlinespace
Workflow & Edit \code{.edp} $\rightarrow$ \code{FreeFem++ waam\_multilayer.edp}
$\rightarrow$ open \code{waam.pvd} in ParaView &
GUI (12 tabs) or CLI $\rightarrow$ slice STL $\rightarrow$ mesh $\rightarrow$
generate $\rightarrow$ solve (serial/MPI) $\rightarrow$ plots + HTML report \\
\addlinespace
Geometry & Hard-coded 26\,$\times$\,3\,mm wall on a
200\,$\times$\,100\,$\times$\,10\,mm plate &
Arbitrary STL/STEP/IGES/BREP, parametric walls, G-code-driven toolpaths \\
\addlinespace
Heat source & Goldak double ellipsoid only & 7 selectable models \\
\addlinespace
Material & Mild steel S235 hard-coded & 8-alloy JSON database, dual
substrate/deposit assignment \\
\addlinespace
Mechanics & Single post-hoc \emph{elastic} solve, inherent strain &
Incremental \textbf{J2 thermo-elasto-plasticity}, Newton with consistent
tangent, TRIP, springback \\
\addlinespace
Thermal extras & Conduction + convection & + radiation, boiling cap,
enthalpy-exact latent heat, adaptive $\Delta t$, Picard correctors,
target-temperature inter-pass control \\
\addlinespace
Beyond original scope & --- & OpenFOAM VOF melt-pool CFD, ultrasonic (UVA)
assistance, microstructure and Vickers hardness prediction \\
\addlinespace
Verification & None & 11-group physics suite, J2 cross-check, external
CAS-DED benchmark, pytest suite \\
\addlinespace
Portability & \code{D:\textbackslash freefem++\textbackslash...} hard-coded &
Binary discovery + environment overrides, MPI, self-contained timestamped run
directories \\
\bottomrule
\end{tabularx}
\end{table}

\noindent\textbf{One-sentence verdict.} The original is a \emph{research
artifact} --- correct, minimal and self-documenting --- while the framework is
an \emph{engineering platform} that retains the original physics as its inner
core and rebuilds everything around it: geometry, toolpath, materials,
plasticity, user interface, verification and reporting.

% ============================================================
\section{The Original: What It Actually Does}
% ============================================================

\subsection{Physics implemented in \file{downloaded.edp}}

\begin{itemize}[leftmargin=1.4em,itemsep=0.15em]
    \item \textbf{3D transient heat equation}, $\theta$-scheme with
    $\theta = 0.5$ (Crank--Nicolson), P1 elements, UMFPACK direct solver
    (lines 196--225).
    \item \textbf{Goldak double-ellipsoid} moving source with
    $a_f = 8$, $a_r = 20$, $b_G = 7$, $c_G = 4$\,mm, $f_f = 0.6$,
    $f_r = 1.4$, a $\tanh$ front/rear split ($\varepsilon_{\text{split}} =
    10^{-3}$) and amplitude
    \begin{equation}
        Q_0 \;=\; \frac{6\sqrt{3}\,f\,P}{\pi\sqrt{\pi}\,a\,b\,c},
        \qquad P = I\,U\,\eta = 200 \times 24 \times 0.8 = 3840\ \text{W}.
        \label{eq:goldak}
    \end{equation}
    \item \textbf{Process schedule}: $v_{\text{scan}} = 10$\,mm/s,
    $t_{\text{pass}} = 20$\,s, $t_{\text{wait}} = 60$\,s,
    $t_{\text{total}} = 2080$\,s, 26 layers.
    \item \textbf{Quiet-element activation}: inactive deposit nodes receive
    $k_{\text{eff}} = 10^{-4}k$ and $C_{p,\text{eff}} = 10^{-4}C_p$; nodes are
    born progressively when $x_{\text{torch}} \ge x_{\text{node}}$ inside the
    weld loop; the source is \emph{masked} by $\psi_{\text{active}}$ so the arc
    cannot heat material that does not yet exist (lines 271--309). The README
    identifies this as the number-one correctness requirement: without it,
    $C_p \approx 0$ produces $10^{7}\,^\circ$C spikes.
    \item \textbf{Temperature-dependent properties}: $k(T)$ falling
    $60 \rightarrow 38 \rightarrow 27$\,W/mK, $C_p(T)$ with a Curie peak of
    5000\,J/kgK at $720\,^\circ$C, $T_{\text{sol}} = 1450\,^\circ$C,
    $T_{\text{liq}} = 1520\,^\circ$C.
    \item \textbf{Boundaries}: $h_{\text{bot}} = 50$ on face~1 (bottom),
    $h_{\text{conv}} = 15$ on faces 2--6. \emph{No radiation, no fusion latent
    heat, no boiling cap} --- conduction and convection only.
    \item \textbf{Fixed three-step timestepping}: $\Delta t_{\text{weld}} =
    0.333$\,s (so $v\Delta t \approx 0.5\,h_x$), $\Delta t_{\text{wait}} =
    10$\,s, $\Delta t_{\text{cool}} = 120$\,s $\times$ 20 steps.
    \item \textbf{Mechanical stage (post-hoc)}: accumulate eigenstrain
    $\varepsilon^{\text{th}} \mathrel{+}= \alpha\,\Delta T$ with
    $\alpha = 2.1\times10^{-5}$ at every step, then a single linear-elastic
    solve with $E(T)$ from 210 to 80\,GPa and $\nu = 0.3$, bottom face fixed,
    written to \code{waam\_mech.vtu} (lines 464--525).
\end{itemize}

\subsection{Mesh and output}

Two \code{cube()} meshes are concatenated with
\code{Th = ThBase + ThDep}: base $30\times20\times8$, deposit
$30\times6\times26$, scaled by linear \code{movemesh3} only (a documented
FreeFEM 4.15 / Windows workaround forbidding trigonometric transforms). The
result is approximately 11{,}687 nodes and 56{,}880 tetrahedra, with one
element per layer in $z$.

Outputs are \code{waam.pvd}, \code{T\_NNNN.vtu} (fields \code{TempC},
\code{Activation}), \code{waam\_mech.vtu} and \code{waam\_data.csv} with eight
columns, all written to the hard-coded path
\code{D:\textbackslash freefem++\textbackslash arc\_waam\_simulation\textbackslash}.

\subsection{Known quirks in the original}
\begin{enumerate}[leftmargin=1.6em,itemsep=0.15em]
    \item \textbf{Header/code mismatch.} The comments promise ``activation via
    a gradual $\tanh$ ramp over 5\,s'' and ``$T = T_{\text{liq}}$ on
    activation''; the code switches $\psi_{\text{active}}$ instantly at
    $x_{\text{torch}} + a_f$ and initialises at $T_{\text{amb}}$
    (deliberately --- see the comment at line 92 warning that
    $T_{\text{dep,init}} = T_{\text{liq}}$ causes double heating).
    \item \textbf{Non-conforming interface.} \code{ThBase + ThDep} merges no
    nodes: the substrate has 5.0\,mm $y$-spacing and the deposit 1.33\,mm, so
    the two blocks are not bonded.
    \item \textbf{No radiation / latent heat / evaporation} despite peak
    temperatures near $1900\,^\circ$C.
    \item \textbf{Elastic-only residual stress} --- no yield limit, so the
    reported von Mises stress can exceed the material's strength.
\end{enumerate}

% ============================================================
\section{Architecture: Monolith to Platform}
% ============================================================

The structural change that dominates every other difference is that
\emph{the physics stopped being the source code and became an emission
target}. \code{EDPBuilder.build()} (\file{codegen/edp\_builder.py}, line 1474)
assembles a script from more than twenty ordered sections:

\begin{quote}\small\ttfamily
GEOMETRY \(\rightarrow\) PROCESS \(\rightarrow\) FEEDSTOCK \(\rightarrow\)
MATERIAL \(\rightarrow\) TIMESTEP \(\rightarrow\) G-CODE READER \(\rightarrow\)
MESH \(\rightarrow\) ADAPTIVE dt \(\rightarrow\) FE SPACES \(\rightarrow\)
\textbf{THERMO-ELASTO-PLASTIC ENGINE} \(\rightarrow\) FORMS \(\rightarrow\)
PROBES \(\rightarrow\) OUTPUT \(\rightarrow\) RESTART \(\rightarrow\)
DEPOSITION LOOP \(\rightarrow\) FINAL COOLING \(\rightarrow\) MECHANICAL
\(\rightarrow\) MICROSTRUCTURE \(\rightarrow\) VALIDATION GATE
\(\rightarrow\) solver\_results.json
\end{quote}

\begin{table}[h!]
\centering
\caption{Largest source files by line count (excluding caches and run
artefacts).}
\label{tab:loc}
\small
\begin{tabular}{@{}r l l@{}}
\toprule
\textbf{Lines} & \textbf{File} & \textbf{Role} \\
\midrule
6\,959 & \file{gui/app.py} & Tkinter desktop GUI, 12 tabs \\
4\,992 & \file{codegen/edp\_builder.py} & FreeFEM script generator \\
1\,958 & \file{core/fluid/openfoam\_coupling.py} & OpenFOAM case generation \\
1\,464 & \file{tests/physics\_suite.py} & 279 automated checks \\
1\,419 & \file{postprocess/plot\_generator.py} & 8 publication plots \\
1\,331 & \file{gui/widgets/waam\_3d\_studio.py} & Three.js WebGL viewer \\
1\,255 & \file{core/postprocess/output\_manager.py} & Summaries, reports \\
1\,242 & \file{gui/widgets/mesh\_viewer\_3d.py} & Embedded PyVista viewport \\
1\,174 & \file{postprocess/report\_generator.py} & HTML engineering report \\
\midrule
\multicolumn{3}{@{}l}{\textit{Original: \file{downloaded.edp} = 534 lines,
\file{README.md} = 434 lines}} \\
\bottomrule
\end{tabular}
\end{table}

Two hand-maintained \code{.edp} files survive as static presets rather than
generator output:
\begin{itemize}[leftmargin=1.4em,itemsep=0.15em]
    \item \file{DEFAULT.edp} / \file{generated\_edps/DEFAULT.edp} (655 lines)
    --- the original's architecture \emph{with fixes}: an enthalpy-exact
    piecewise-linear $C_p$ band over $[T_{\text{sol}}, T_{\text{liq}}]$, a
    quiet method that quenches only $k$ (retaining full mass so diffusivity
    collapses instead of thermal mass), linearised Stefan--Boltzmann radiation
    ($\varepsilon = 0.7$), a boiling cap at $T_{\text{boil}} = 2862\,^\circ$C
    with an $E_{\text{cap}}$ energy audit, and an explicit caveat that the
    mechanical stage is a ``post-hoc ELASTIC estimate (no plasticity)'';
    \item \file{generated\_edps/uva\_waam\_ss308l\_ss316.edp} (859 lines) ---
    SS316 substrate with SS308L deposit over 33 layers and \emph{ultrasonic
    vibration assistance} (20\,kHz, 15\,$\mu$m, Blaha yield softening).
\end{itemize}

% ============================================================
\section{Geometry and Meshing}
% ============================================================

\begin{table}[h!]
\centering
\caption{Geometry and mesh pipeline comparison.}
\label{tab:mesh}
\small
\begin{tabularx}{\linewidth}{@{}l X X@{}}
\toprule
& \textbf{Original} & \textbf{Framework} \\
\midrule
Input & None (hard-coded) & STL, STEP, STP, IGES, IGS, BREP, PLY, OBJ, G-code \\
\addlinespace
Mesher & FreeFEM \code{cube()} $\times 2$ & Gmsh OCCT volume meshing
(\file{mesh/stl\_pipeline.py:541}) \emph{or} pure-Python graded structured
mesh (\file{mesh/structured\_conforming.py:464}) \\
\addlinespace
Sizing & Fixed $30\times20\times8$ / $30\times6\times26$ & Rule-based engine,
3 strategies, node budget 1500/2500/5000/custom
(\file{mesh/intelligent\_mesher.py:29}) \\
\addlinespace
Interface & \code{ThBase + ThDep} $\rightarrow$ \textbf{zero-width crack} &
\textbf{Single conforming mesh, nodes shared by construction} \\
\addlinespace
Labels & Implicit & Enforced: 1 = bottom, 5 = top, 2/3/4/6 = flanks, shared by
mesher and generated forms \\
\addlinespace
Quality control & None & Signed-volume, orientation and outward-normal checks;
the mesher \emph{raises} rather than emit an unusable mesh
(\file{edp\_builder.py:3168}) \\
\bottomrule
\end{tabularx}
\end{table}

The conforming-mesh work is the framework's most consequential \emph{correctness}
fix. The crack produced by the original's mesh union was measured on a running
model at a \textbf{19.9\,$^\circ$C mean and 92.9\,$^\circ$C peak} temperature
discontinuity across the interface --- i.e.\ no genuine conduction and no
mechanical bond between substrate and deposit.

% ============================================================
\section{Toolpath: Linear Scan to Slicer and G-code}
% ============================================================

\begin{itemize}[leftmargin=1.4em,itemsep=0.2em]
    \item \textbf{Original:} the torch traverses $x = 0 \rightarrow 200$\,mm in
    one direction, every layer, at constant $v_{\text{scan}}$.
    \item \textbf{Framework:} \file{mesh/slicer\_api.py} (695 lines) slices
    geometry into layers using five infill patterns (zigzag, meander, contour,
    crosshatch, unidirectional) with overlap, perimeters and bidirectional
    reversal, emitting \code{toolpath.gcode}, \code{toolpath.txt} and
    \code{layer\_times.txt}. \file{mesh/gcode\_parser.py} (411 lines) parses
    G0/G1/\textbf{G2/G3 arcs} into solver arrays.
    \item The generated script reads \code{toolpath.txt} ($t$, $x$, $y$, $z$,
    power) and \code{layer\_pass\_times.txt}, gates the arc off during rapid
    travels, supports a direction sign for zig-zag layers, and clamps the
    activation lookup with an ``afterProgram'' guard so an over-energised arc
    is not re-applied after the program ends.
\end{itemize}

% ============================================================
\section{Thermal Physics}
% ============================================================

\begin{table}[h!]
\centering
\caption{Thermal formulation comparison.}
\label{tab:thermal}
\small
\begin{tabularx}{\linewidth}{@{}l X X@{}}
\toprule
\textbf{Aspect} & \textbf{Original} & \textbf{Framework} \\
\midrule
Scheme & $\theta = 0.5$, one solve per step & $\theta$-scheme + Picard
correctors (3 passes, $0.5\,^\circ$C tolerance) for non-linearity \\
\addlinespace
Timestepping & Three fixed steps (0.333 / 10 / 120\,s) & Fixed \emph{or}
adaptive Fourier-number controller (Fo target 0.45, torch-advance cap,
runaway guard) \\
\addlinespace
Latent heat & $C_p$ spike to 5000\,J/kgK at the Curie point only; no fusion
latent & Enthalpy-exact piecewise-linear $C_p$ band over
$[T_{\text{sol}}, T_{\text{liq}}]$ preserving $\int C_p\,dT$
(fixes $\sim$10\,\% energy error) \\
\addlinespace
Radiation & None & Linearised Stefan--Boltzmann on all outer surfaces,
emissivity per alloy \\
\addlinespace
Evaporation & None & Boiling cap at $T_{\text{boil}}$ with $E_{\text{cap}}$
energy accounting \\
\addlinespace
Inter-pass control & Fixed $t_{\text{wait}} = 60$\,s & \code{time} mode (as
original) \emph{or} \code{target} mode --- iterative cooldown to a target peak
temperature (default $400\,^\circ$C) with convergence test and iteration cap \\
\addlinespace
Quiet elements & $k$ and $C_p$ both scaled by $10^{-4}$ & $k$ quenched with
full mass retained; plus optional prescribed-$T$ birth (Ansys-style liquidus
gate) \\
\addlinespace
Energy audit & \code{Tmax} printed to the console & \code{solver\_results.json}
balances intended arc, injected arc, stored enthalpy, boiling-cap removal,
surface losses, quiet-node discard, floor additions and numerical residual \\
\bottomrule
\end{tabularx}
\end{table}

% ============================================================
\section{Heat Sources}
% ============================================================

\begin{table}[h!]
\centering
\caption{Available volumetric and surface heat-source models.}
\label{tab:src}
\small
\begin{tabularx}{\linewidth}{@{}l c c X@{}}
\toprule
\textbf{Model} & \textbf{Orig.} & \textbf{Fwk.} & \textbf{Notes} \\
\midrule
Goldak double ellipsoid & \yes & \yes & Identical algebra to
Eq.~\eqref{eq:goldak} --- direct carry-over \\
Gaussian volumetric & \no & \yes & $Q_0 = 3\sqrt{3}P/(\pi\sqrt{\pi}R^3)$ \\
Gaussian surface & \no & \yes & Neumann \code{int2d} boundary condition \\
Conical frustum (keyhole) & \no & \yes & $R(\zeta) = r_t + (r_b - r_t)\zeta/d_{\text{pen}}$ \\
Uniform band & \no & \yes & $\tanh$-smoothed box \\
Prescribed temperature & \no & \yes & Post-solve $T = T_{\text{liq}}$ inside the activation radius \\
P48 arc/droplet partition & \no & \yes & \emph{Prog.\ Addit.\ Manuf.}\ 1 (2025) 10:11647; empirical $\sigma_q = 6.205\times10^{-4} I^{0.28}$ \\
\bottomrule
\end{tabularx}
\end{table}

\noindent All models remain masked by $\psi_{\text{active}}$ --- the original's
key insight, preserved intact. Automatic source sizing keeps the semi-axes at or
above the local element size with a fidelity floor
(\code{\_auto\_size\_source}, \file{edp\_builder.py:468}). The six files under
\code{core/heat\_sources/*.edp} are explicitly marked \emph{standalone reference
modules --- not wired into the pipeline}; the live physics is whatever
\code{EDPBuilder} emits.

% ============================================================
\section{Materials}
% ============================================================

\begin{table}[h!]
\centering
\caption{Material data comparison.}
\label{tab:mat}
\small
\begin{tabularx}{\linewidth}{@{}l X X@{}}
\toprule
& \textbf{Original} & \textbf{Framework} \\
\midrule
Source & \code{kSteel(T)}, \code{CpSteel(T)} written inline &
\file{core/multimaterial/material\_database.json} (1{,}309 lines) \\
\addlinespace
Alloys & 1 (S235/A36 mild steel) & 8: AlSi10Mg, ER70S-6, IN625, Pure Al,
SS304, SS308L, SS316L, Ti-6Al-4V \\
\addlinespace
Data & $k$, $C_p$, $\rho$, $T_{\text{sol}}$, $T_{\text{liq}}$, $E$, $\nu$ &
+ $\nu$, latent heat, emissivity, composition, temperature tables for
$k$/$E$/$\sigma_y$/$\alpha$, phase transformation, hardness coefficients, wire
diameters \\
\addlinespace
Assignment & One material for the whole domain & \textbf{Dual}: substrate and
deposit alloys chosen independently, with a dissimilar-metal
CTE/thermophysical compatibility report \\
\addlinespace
Phase kinetics & None & \file{core/multimaterial/alloy\_metallurgy.py}: 5
phase families $\times$ 4 phases (Burgers BOR, WRC-1992 + Hall-Petch,
Scheil/Laves, cellular eutectic, Kirkaldy-Venugopalan + Koistinen-Marburger),
IIW carbon equivalent, WRC Cr/Ni equivalents \\
\bottomrule
\end{tabularx}
\end{table}

% ============================================================
\section{Mechanics: Elastic Estimate to Incremental Plasticity}
% ============================================================

\begin{table}[h!]
\centering
\caption{Mechanical formulation comparison.}
\label{tab:mech}
\small
\begin{tabularx}{\linewidth}{@{}l X X@{}}
\toprule
& \textbf{Original} & \textbf{Framework} \\
\midrule
Formulation & One post-hoc linear-elastic solve &
Incremental J2 thermo-elasto-plasticity, backward-Euler radial return \\
\addlinespace
Tangent & n/a & Consistent-tangent Newton (default) or Picard, with force and
displacement residuals recorded \\
\addlinespace
Hardening & None & Isotropic $H'$ \\
\addlinespace
TRIP & None & Transformation strain $\beta = 0.011$ \\
\addlinespace
Remelt & Eigenstrain keeps growing unphysically & Plastic memory reset on
remelt; above-solidus ``slurry'' stiffness cut avoids singular degrees of
freedom \\
\addlinespace
Eigenstrain & $\varepsilon^{\text{th}}$ accumulated from $t=0$, can exceed
yield $\rightarrow$ phantom stress & Correct eigenstress subtraction;
\file{DEFAULT.edp} documents the phantom-stress fix \\
\addlinespace
Boundary conditions & Bottom fixed only & Isostatic / fixture-clamped /
cantilever, fixture stiffness, springback with fixture-release regularisation,
gravity \\
\addlinespace
Discretisation & P1 only & P1 or P2; legacy inherent-strain fallback retained
but disabled \\
\addlinespace
Output & \code{waam\_mech.vtu}, maxima to console & + PEEQ, per-frame von
Mises, Newton convergence in \code{solver\_results.json}, checkpoint/restart \\
\bottomrule
\end{tabularx}
\end{table}

% ============================================================
\section{New Physics Domains Absent from the Original}
% ============================================================

\begin{enumerate}[leftmargin=1.8em,itemsep=0.35em]
    \item \textbf{OpenFOAM melt-pool CFD.} \file{core/fluid/openfoam\_coupling.py}
    (1{,}958 lines) generates an OpenFOAM 2412 \code{beamWeldFoam} case:
    volume-of-fluid two-phase (metal plus atmosphere),
    diverging-oscillating-Gaussian melting source, Marangoni shear, vapour
    recoil and a Darcy mushy-zone sink. It maps
    $I\,U\,\eta \rightarrow$ \code{MELTING/HS\_Q},
    $v_{\text{scan}} \rightarrow$ \code{HS\_velocity} and
    $b_G \rightarrow$ \code{HS\_a}, then writes \code{cfd\_history.csv},
    plots and \code{cfd\_report.html}.
    \item \textbf{Ultrasonic vibration assistance (UVA).} 20\,kHz /
    15\,$\mu$m with four acoustic modules (streaming, heat generation,
    cavitation, momentum transfer); amplitude acts on the \emph{yield surface}
    (Blaha softening) rather than on Young's modulus, with a layer-by-layer
    multi-material zone table.
    \item \textbf{Microstructure and hardness.} G/R solidification map, JMAK
    kinetics, secondary dendrite-arm-spacing laws per alloy family, reheat
    anneal factors and discrete Vickers indentation predictions, producing
    \code{hardness\_profile\_z.csv} and \code{hardness\_test\_locations.csv}.
    \item \textbf{Reference modules not yet wired:} Marangoni melt pool,
    vibration acoustics and solidification-phase \code{.edp} files under
    \code{core/}; the \code{core/thermal/} and \code{core/mechanical/}
    directories are empty.
\end{enumerate}

% ============================================================
\section{Output and Postprocessing}
% ============================================================

\begin{table}[h!]
\centering
\caption{Output artefacts and postprocessing.}
\label{tab:out}
\small
\begin{tabularx}{\linewidth}{@{}l X X@{}}
\toprule
& \textbf{Original} & \textbf{Framework} \\
\midrule
Fields & \code{TempC}, \code{Activation}, \code{Ux Uy Uz DispMag SigMises} &
+ PEEQ, per-frame mechanics, UVA fields, \code{VickersHV} \\
\addlinespace
Files & \code{waam.pvd}, \code{T\_*\_*.vtu}, \code{waam\_mech.vtu},
\code{waam\_data.csv} & + \code{solver\_results.json},
\code{run\_summary.json}, \code{parameters.json}, \code{probe\_points.csv},
\code{probe\_paths.csv}, hex-VTK conversions, \code{waam\_3d\_studio.html} \\
\addlinespace
Plots & None --- the README asks the user to write matplotlib by hand &
8 publication figures at 300\,dpi: thermal history and peaks, cooling rates
and $t_{8/5}$, distortion and warpage, von Mises, heat input and dwell energy,
baseplate thermal cycle with elevation probes, spatial stress paths (sections
A--A, B--B), normalised warpage index \\
\addlinespace
Reports & None & \code{report.html} (two-page A4), \code{run\_summary.json},
ParaView loader script, \code{cfd\_report.html}, toolpath GIF, uncut 3D
renders \\
\addlinespace
Visualisation & Manual ParaView & Embedded PyVista viewport, standalone
PyVista GPU window, Three.js WebGL studio (assets inlined for offline use) with
heat-source inspector and G-code playback $1\times$--$50\times$; ParaView still
one click away \\
\addlinespace
History output & None & Abaqus-style \code{*OUTPUT,HISTORY} equivalent:
probe points and paths sampled by the solver at true simulated times \\
\bottomrule
\end{tabularx}
\end{table}

% ============================================================
\section{Workflow and User Experience}
% ============================================================

\begin{table}[h!]
\centering
\caption{Workflow comparison.}
\label{tab:ux}
\small
\begin{tabularx}{\linewidth}{@{}l X X@{}}
\toprule
& \textbf{Original} & \textbf{Framework} \\
\midrule
Entry point & Edit the file, run one command & \code{run.py} with ten modes:
default (GUI), \code{--cli}, \code{--stl}, \code{--slice}, \code{--validate},
\code{--validate-full}, \code{--generate-only}, \code{--openfoam},
\code{--uva}, \code{--version} \\
\addlinespace
Parameters & 40+ constants at the top of the file & Live GUI with 8 presets,
units-aware mm/SI conversion, \code{.waam} project save and load \\
\addlinespace
GUI & None & Tkinter, 12 tabs: Geometry, Process, Material, Mesh/CAD, Physics,
Solver, OpenFOAM CFD, Parametric and DOE, 3D Viewport, Script Preview,
Ultrasonic, Simulation Log --- plus live console, progress parsing, themes and
tooltips \\
\addlinespace
Orchestration & Manual & Five stages --- slice $\rightarrow$ mesh $\rightarrow$
generate $\rightarrow$ solve $\rightarrow$ postprocess --- into a
self-contained \code{waam\_output/<timestamp>/} with a \code{latest} pointer \\
\addlinespace
Persistence & Rewrite the file & \code{.waam} projects,
\code{parameters.json} as single source of truth, checkpoint/restart by layer
or every $N$ steps \\
\addlinespace
Portability & Hard-coded \code{D:\textbackslash freefem++} &
\code{find\_freefem\_serial/mpi()} plus \code{FREEFEM\_BIN},
\code{FF\_MPIRUN}, \code{PARAVIEW\_BIN} overrides; macOS bundle default; a
pseudo-terminal wrapper so FreeFEM's standard output flushes line by line \\
\bottomrule
\end{tabularx}
\end{table}

\noindent A telling detail of continuity: the GUI's progress bar parses
\emph{the original script's log lines} (\code{--- Layer n/N},
\code{ws=... Tmax=...}, \code{Inter-pass}), \file{gui/app.py:6112--6151}.

% ============================================================
\section{Verification and Validation}
% ============================================================

\begin{table}[h!]
\centering
\caption{Verification suites present in the framework (none exist upstream).}
\label{tab:vv}
\small
\begin{tabularx}{\linewidth}{@{}l X@{}}
\toprule
\textbf{Suite} & \textbf{What it checks} \\
\midrule
\file{tests/physics\_suite.py} (1{,}464 lines) & 11 groups, 279 assertions:
heat-source normalisation, latent enthalpy, VOF, ferrite number, builder
markers, static-preset integrity, \code{.edp} module structure, incremental
mechanics, history/restart, structured results, conforming interface \\
\addlinespace
\file{tests/validate\_j2\_mechanics.py} (203) & Elastic baseline against J2
plasticity (radial return, TRIP, annealing) on an identical two-layer mesh \\
\addlinespace
\file{tests/validate\_cas\_ded\_dimeen26.py} (599) & \textbf{External
benchmark CAS-DEDD-IMEEN26} (Ansys / TU Wien 2026, AlSi10Mg DED): reference
targets 0.25\,mm max deformation, $4\times10^{-2}$\,mm mean, 580\,$^\circ$C
process temperature, 50\,mm$^3$/s, 2.5\,mm bead radius. Executes a real FreeFEM
run, therefore gated behind \code{--validate-full} \\
\addlinespace
\file{tests/test\_wall\_probes\_ultrasonic.py} (219) & pytest: deposition
-pattern codegen, ultrasonic gating, probe stations, plot generation \\
\addlinespace
\file{parametric\_sweep.py} (324) & Six cases varying current, voltage, wire
feed, speed, bead dimensions and layer count, with an incrementally written
\code{sweep\_timing\_summary.json} \\
\bottomrule
\end{tabularx}
\end{table}

% ============================================================
\section{Lineage: What Came From the Original}
% ============================================================

\subsection{Direct carry-over}

\begin{enumerate}[leftmargin=1.8em,itemsep=0.2em]
    \item The $\theta$-scheme 3D transient variational forms --- same structure
    as \code{aThermal}/\code{lThermal}, \file{downloaded.edp:199--225}.
    \item The Goldak double ellipsoid with $\tanh$ front/rear split and
    $f_f = 0.6$, $f_r = 1.4$ --- identical algebra.
    \item Quiet-element activation with $\psi_{\text{active}}$ source masking
    --- the original's central correctness insight, retained everywhere.
    \item Progressive $x$-based node activation inside the weld loop.
    \item The inter-pass cooling cadence and the weld / wait / cooldown stage
    structure of the deposition loop.
    \item Cumulative thermal eigenstrain accumulation
    $\varepsilon^{\text{th}} \mathrel{+}= \alpha\,\Delta T$.
    \item The boundary-label scheme (1 bottom, 5 top, 2/3/4/6 sides) and the
    PVD / VTU / \code{waam\_data.csv} output contract.
    \item The console progress format --- which the GUI now parses.
    \item The default geometry: $n_{\text{layers}} = 26$, 200\,mm length,
    8\,mm bead width, 200\,A / 24\,V / $\eta = 0.8$ / 10\,mm/s.
\end{enumerate}

\subsection{Evolved or corrected}

\begin{table}[h!]
\centering
\caption{Places where the framework deliberately departs from the original.}
\label{tab:fixes}
\small
\begin{tabularx}{\linewidth}{@{}X X@{}}
\toprule
\textbf{Original} & \textbf{Framework} \\
\midrule
$C_p$ and $k$ both scaled by $10^{-4}$ in quiet elements & Quench only $k$,
keep full mass: diffusivity collapses, thermal mass does not \\
\addlinespace
No fusion latent heat & Enthalpy-exact piecewise-linear $C_p$ band over
$[T_{\text{sol}}, T_{\text{liq}}]$ \\
\addlinespace
No radiation or boiling cap & Stefan--Boltzmann plus $T_{\text{boil}}$ clip
with $E_{\text{cap}}$ audit \\
\addlinespace
Eigenstrain may exceed yield $\rightarrow$ phantom stress & Correct
eigenstress subtraction, or full J2 plasticity \\
\addlinespace
\code{ThBase + ThDep} non-conforming interface & Single conforming mesh \\
\addlinespace
Hard-coded \code{D:\textbackslash} output & Portable, timestamped run
directories \\
\bottomrule
\end{tabularx}
\end{table}

% ============================================================
\section{Gap Analysis: What the Framework Still Lacks}
% ============================================================

\begin{enumerate}[leftmargin=1.8em,itemsep=0.3em]
    \item \textbf{No top-level README} comparable to the original's, which
    documents physics, geometry, ParaView recipes, an error table, extension
    ideas and a citation. That knowledge now sits inside GUI dialogs and code
    comments.
    \item \textbf{\code{core/heat\_sources/*.edp} and
    \code{core/\{thermal,mechanical\}} are unwired} --- the reference modules
    are never executed by the pipeline.
    \item \textbf{\code{waam\_parametric\_sweep\_comparison\_report.pdf} is not
    reproducible} --- no code in the repository generates it.
    \item \textbf{Two millimetre-detection heuristics coexist} ---
    \file{mesh/units.py} per-key thresholds versus the blanket
    \code{>0.05} / \code{>0.1} test in \file{run.py}.
    \item \textbf{\code{test\_wall\_probes\_ultrasonic.py} is outside
    \code{--validate}} --- a routine health check misses it unless pytest is
    run separately.
    \item \textbf{Multi-material per-layer alloy switching is declared but
    unimplemented}; only the zonal substrate/deposit split works.
    \item \textbf{The report is HTML, not PDF}, despite ``LaTeX PDF-style''
    styling.
    \item \textbf{Licensing} --- see Section~\ref{sec:license}.
\end{enumerate}

% ============================================================
\section{Licensing and Provenance}
\label{sec:license}
% ============================================================

\begin{table}[h!]
\centering
\caption{Licensing of the two works.}
\label{tab:lic}
\small
\begin{tabularx}{\linewidth}{@{}l X X@{}}
\toprule
& \textbf{Original} & \textbf{Framework} \\
\midrule
License & CC BY-NC 4.0 & GNU GPL v3.0 or later \\
\addlinespace
Holder & Akshansh Mishra (akshansh11) & Pranay1004 and the WAAM\_IIST
contributors \\
\addlinespace
DOI / URL & \url{https://doi.org/10.5281/zenodo.20753459} &
\url{https://github.com/Pranay1004/WAAM_IIST} \\
\addlinespace
Commercial use & Not permitted (NC field) & Permitted; derivatives must remain
GPL \\
\addlinespace
Files in this repository & \code{original\_git\_file/} retained verbatim &
\code{license/LICENSE}, \code{license/README.md}, \code{license/NOTICE} \\
\bottomrule
\end{tabularx}
\end{table}

\noindent CC BY-NC 4.0 carries a \emph{NonCommercial} field, while the GPL
deliberately forbids additional restrictions; the two are therefore not
license-compatible in the direction of relicensing upstream code as GPL. The
repository does \emph{not} relicense upstream: \code{original\_git\_file/} is
kept byte-for-byte under CC BY-NC 4.0, only the framework's own source is
offered under GPL-3.0-or-later, and a recipient must honour both regimes. The
algorithmic ideas themselves (Goldak source, quiet-element activation,
progressive torch-triggered birth) are uncopyrightable \emph{methods} and are
reimplemented from scratch in \file{codegen/edp\_builder.py}; no upstream source
text is copied into the generator. Third-party components keep their own terms
(three.js and OrbitControls: MIT; FreeFEM++ and OpenFOAM: GPL, invoked
externally and never redistributed).

% ============================================================
\section{Verdict}
% ============================================================

\begin{table}[h!]
\centering
\caption{Qualitative scoring of the two code bases.}
\label{tab:verdict}
\small
\begin{tabularx}{\linewidth}{@{}l c c X@{}}
\toprule
\textbf{Dimension} & \textbf{Orig.} & \textbf{Fwk.} & \textbf{Comment} \\
\midrule
Correctness of the core thermal solve & Good & Stronger &
Conforming mesh, energy audit, latent heat, radiation close the original's four
physical omissions \\
\addlinespace
Fidelity of the mechanics & Weak & Strong & Elastic estimate versus validated
incremental J2 \\
\addlinespace
Breadth of physics & Narrow & Broad & + CFD, UVA, microstructure, hardness \\
\addlinespace
Realism of geometry and toolpath & Poor & Strong & CAD/G-code, five infill
patterns \\
\addlinespace
Usability & Weak & Strong & GUI, reports, 3D viewers, project files \\
\addlinespace
Verification & None & Strong & Four suites plus an external benchmark \\
\addlinespace
Documentation & Excellent & Fair & The original's single README was not
carried forward \\
\addlinespace
Accessibility of the physics & Excellent & Weak & 534 readable lines versus a
4{,}992-line generator \\
\addlinespace
Dependency weight & Minimal & Heavy & FreeFEM++ only, versus FreeFEM + Gmsh +
PyVista + VTK + matplotlib + three.js (+ optional OpenFOAM, pywebview) \\
\bottomrule
\end{tabularx}
\end{table}

\noindent\textbf{Conclusion.} The framework is a strict superset of the
original: it keeps the original's solver engine (FreeFEM++), Goldak and
quiet-element physics, output contract and even its log format, then adds
everything a usable simulation platform requires. Its one regression relative
to the original is \emph{legibility} --- 534 self-contained, heavily commented
lines became a 4{,}992-line generator plus forty Python modules, and the
original's exemplary top-level README was not carried forward. The remedy
suggested by this comparison is modest: restore a single narrative README at
the repository root that states the physics, the file map and the validation
status in the same voice as the upstream document.

% ============================================================
\section*{References}
% ============================================================
\begin{enumerate}[leftmargin=2.2em,itemsep=0.25em,label={[\arabic*]}]
    \item A.~Mishra. \emph{WAAM-MultiLayer-Simulation}. Zenodo, 2026.
    \url{https://doi.org/10.5281/zenodo.20753459}. Licensed CC BY-NC 4.0.
    \item P.~K.~Pandey and the WAAM\_IIST contributors. \emph{WAAM-Framework /
    WAAM\_IIST}, 2026. \url{https://github.com/Pranay1004/WAAM_IIST}.
    Licensed GPL-3.0-or-later.
    \item M.~Goldak, A.~Chakravarti and M.~Bibby. A new finite element model
    for welding heat sources. \emph{Metallurgical Transactions B},
    17B:187--193, 1986.
    \item Progress in Additive Manufacturing, 1:10:11647, 2025 (P48 arc/droplet
    energy-partition heat source).
    \item CAS-DEDD-IMEEN26 external validation case (Ansys Academic Program /
    TU Wien, 2026, AlSi10Mg DED).
    \item FreeFEM++ finite-element solver, \url{https://freefem.org},
    version 4.14/4.15.
    \item OpenFOAM \code{beamWeldFoam}, OpenCFD Ltd., version 2412.
\end{enumerate}

\end{document}
